% answers_short.tex —— 1992 年·上海卷 · 简短答案速查
% 由 tools/make_short_answers.py 自 answers.tex 的 \gksolhead/\gkanswer 提取，
% 并与 1992/1992_上海_答案要点.md 逐题比对后生成，勿手改（改源 MD 后重跑）。

\begin{gknotice}
\textbf{使用说明}：本节只给最终答案便于快速核对，完整推导见解析版。答案为固定来源 PDF 卷内「解答及评分标准」所载（出版者/版次未核实），非官方标准答案；第 12 题选项指数按卷内答案判读。
\end{gknotice}

\gkpart{填空题}
\begin{itemize}[leftmargin=2.6em]
  \item[\textbf{1.}] $-\dfrac{\sqrt{2}}{4}$。
  \item[\textbf{2.}] $56$。
  \item[\textbf{3.}] $\dfrac{3}{2}$。
  \item[\textbf{4.}] $4$。
  \item[\textbf{5.}] $-64$。
  \item[\textbf{6.}] $x-y+1=0$。
  \item[\textbf{7.}] $60\sqrt{2}\,\pi$。
  \item[\textbf{8.}] $[-1,+\infty)$。
  \item[\textbf{9.}] $72$。
  \item[\textbf{10.}] $\arcsin\dfrac{\sqrt{6}}{4}$。
\end{itemize}

\gkpart{选择题}
\begin{center}\begin{tabular}{ccccccccccc}
\toprule
题号 & 11 & 12 & 13 & 14 & 15 & 16 & 17 & 18 & 19 & 20 \\
\midrule
答案 & D & C & D & D & C & A & B & A & B & B \\
\bottomrule
\end{tabular}\end{center}

\gkpart{解答题}
\begin{itemize}[leftmargin=2.6em]
  \item[\textbf{21.}] $\alpha+\beta=\pi+\arccos\dfrac{33}{65}$。
  \item[\textbf{22.}] $\theta\in\left[0,\dfrac{\pi}{4}\right]\cup\left[\dfrac{7\pi}{4},2\pi\right)$。
  \item[\textbf{23.}] 轨迹为椭圆 $\dfrac{x^{2}}{2}+y^{2}=1$，以及椭圆 $\dfrac{x^{2}}{6}+\dfrac{y^{2}}{3}=1$ 夹在直线 $x=2$、$x=-2$ 之间的部分。
  \item[\textbf{24.}] （1）见解答中的证明；（2）$\dfrac{8}{3}\pi$；（3）$60^{\circ}$。
  \item[\textbf{25.}] （1）当 $q=1$ 时，$\lim\limits_{n\to\infty}\dfrac{S_{n}}{G_{n}}=\dfrac{1}{a}$；当 $0<q<1$ 时，$\lim\limits_{n\to\infty}\dfrac{S_{n}}{G_{n}}=\dfrac{1+q}{a}$；（2）当 $n=1$ 及 $n\ge 4$（$n\in N$）时，$S_{n}<3na_{n}$；当 $n=2,3$ 时，$S_{n}>3na_{n}$。
  \item[\textbf{26.}] $\dfrac{\sqrt{5}-1}{4}$。
\end{itemize}
